% !TeX program = pdflatex
% 中央财经大学 Beamer 模板（非官方）
% 版式参考：Zhuming Shi 的 PKU-Beamer-Template（CC BY 4.0）。
% https://www.overleaf.com/latex/templates/bei-da-zhong-wen-mo-ban-pku-beamer-template/kfxpbtzrqhrn
% 原模板来源：Federico Zenith 的 SINTEF Presentation；Liu Qilong 的衍生主题。
% 本版独立实现主题；校徽与校名图片位于 assets 文件夹。
% 校徽来源：https://www.cufe.edu.cn/info/1029/1066.htm
% 默认使用 pdfLaTeX；中文通过 CJKutf8 与 gbsn 字体排版。
\documentclass[aspectratio=169,11pt]{beamer}
\usepackage[T1]{fontenc}
\usepackage{lmodern}
\usepackage{CJKutf8}
\usepackage{amsmath,amssymb,booktabs,graphicx,tikz}
\usetikzlibrary{positioning,arrows.meta}
\usefonttheme[onlymath]{serif}

% 设计用屏幕配色，可在这里统一调整；非官方标准色声明。
\definecolor{cufeblue}{HTML}{123F68}
\definecolor{cufedark}{HTML}{092C4C}
\definecolor{cufelight}{HTML}{EDF3F8}
\definecolor{cufegold}{HTML}{B78A46}
\definecolor{cufetext}{HTML}{263442}
\definecolor{cufegray}{HTML}{657587}
\colorlet{maincolor}{cufeblue}
\setbeamercolor{normal text}{fg=cufetext,bg=white}
\setbeamercolor{structure}{fg=cufeblue}
\setbeamercolor{alerted text}{fg=cufegold}
\setbeamercolor{frametitle}{fg=cufeblue,bg=white}
\setbeamercolor{framesubtitle}{fg=cufegray}
\setbeamercolor{block title}{fg=white,bg=cufeblue}
\setbeamercolor{block body}{fg=cufetext,bg=cufelight}
\setbeamercolor{block title alerted}{fg=white,bg=cufegold}
\setbeamercolor{block body alerted}{fg=cufetext,bg=cufegold!10}
\setbeamercolor{block title example}{fg=white,bg=cufedark}
\setbeamercolor{block body example}{fg=cufetext,bg=cufelight}
\setbeamerfont{title}{size=\LARGE,series=\bfseries}
\setbeamerfont{subtitle}{size=\large}
\setbeamerfont{frametitle}{size=\Large,series=\bfseries}
\setbeamerfont{framesubtitle}{size=\small}
\setbeamertemplate{navigation symbols}{}
\setbeamertemplate{blocks}[default]
\setbeamertemplate{itemize item}{\raisebox{.15ex}{\scriptsize$\blacktriangleright$}}
\setbeamertemplate{itemize subitem}{\scriptsize$\bullet$}
\setbeamertemplate{enumerate items}[default]
\setbeamersize{text margin left=8mm,text margin right=8mm}
\hypersetup{colorlinks=true,urlcolor=cufeblue,linkcolor=cufeblue}

% 校徽与书法校名使用透明背景 PNG；校名通过 trim 去除多余透明留白。
\newcommand{\cufeseal}[1]{\includegraphics[width=#1]{assets/cufe-logo-transparent.png}}
\newcommand{\cufewordmark}[1]{\includegraphics[width=#1,trim=0 154bp 0 148bp,clip]{assets/cufe-wordmark-transparent.png}}

\setbeamertemplate{frametitle}{%
  \vspace{2mm}
  \begin{beamercolorbox}[wd=\textwidth]{frametitle}
    \usebeamerfont{frametitle}\insertframetitle\strut\par
    \ifx\insertframesubtitle\empty\else
      {\usebeamerfont{framesubtitle}\color{cufegray}\insertframesubtitle\par}
    \fi
    \vspace{1.2mm}{\color{cufeblue}\rule{\textwidth}{.6pt}}
  \end{beamercolorbox}
}
\newcommand{\cufefootertext}{中央财经大学}
\newcommand{\footlinepayoff}[1]{\renewcommand{\cufefootertext}{#1}}
\newcommand{\footlinecolor}[1]{%
  \if\relax\detokenize{#1}\relax
    \setbeamercolor{cufefoot}{fg=cufegray,bg=white}%
  \else
    \setbeamercolor{cufefoot}{fg=white,bg=#1}%
  \fi
}
\footlinecolor{cufeblue}
\setbeamertemplate{footline}{%
  % Beamer 在正文中文环境启动前测量页脚，需在此独立启用中文。
  \begin{CJK*}{UTF8}{gbsn}
  \begin{beamercolorbox}[wd=\paperwidth,ht=2.7mm,dp=1.4mm,leftskip=8mm,rightskip=8mm]{cufefoot}
    \fontsize{6.5}{8}\selectfont
    \cufefootertext\hspace{3mm}|\hspace{3mm}\insertshorttitle
    \hfill\insertframenumber\,/\,\inserttotalframenumber
  \end{beamercolorbox}
  \end{CJK*}
}

% 可选校园图片；不设置时使用纯蓝色斜切面。
\newcommand{\cufecoverimage}{}
\newcommand{\cufesetbackground}[1]{\renewcommand{\cufecoverimage}{#1}}
% 兼容原模板的带星号封面图片接口。
\makeatletter
\newcommand{\titlebackground}{\@ifstar\cufesetbackground\cufesetbackground}
\makeatother

\setbeamertemplate{title page}{%
  \begin{tikzpicture}[remember picture,overlay,x=\paperwidth,y=\paperheight]
    \begin{scope}[shift={(current page.south west)}]
      \fill[cufeblue] (.59,0)--(1,0)--(1,1)--(.77,1)--cycle;
      \ifx\cufecoverimage\empty\else
        \begin{scope}
          \clip (.59,0)--(1,0)--(1,1)--(.77,1)--cycle;
          \node[anchor=center,inner sep=0] at (.81,.5)
            {\includegraphics[height=\paperheight]{\cufecoverimage}};
          \fill[cufeblue,opacity=.55] (.59,0) rectangle (1,1);
        \end{scope}
      \fi
      \node[anchor=north west,inner sep=0] at (.065,.94)
        {\cufeseal{.95cm}};
      \node[anchor=west,inner sep=0]
        at (.145,.895) {\cufewordmark{.28\paperwidth}};
      \node[anchor=west,text=cufegray,inner sep=0,font=\fontsize{5.7}{7}\selectfont]
        at (.145,.820) {CENTRAL UNIVERSITY OF FINANCE AND ECONOMICS};
      \node[anchor=west,inner sep=0] at (.065,.49) {%
        \begin{minipage}{.48\paperwidth}
          {\color{cufeblue}\usebeamerfont{title}\inserttitle\par}
          \vspace{2mm}
          {\color{cufeblue}\usebeamerfont{subtitle}\insertsubtitle\par}
          \vspace{5mm}
          {\small\bfseries\insertauthor\par}
          \vspace{2mm}
          {\small\insertinstitute\par}
          \vspace{2mm}
          {\small\insertdate\par}
        \end{minipage}};
      \node[anchor=south west,text=cufegray,font=\fontsize{7}{9}\selectfont,inner sep=0]
        at (.065,.06) {忠诚\quad 团结\quad 求实\quad 创新};
    \end{scope}
  \end{tikzpicture}
}
\renewcommand{\maketitle}{\begin{frame}[plain,noframenumbering]\titlepage\end{frame}}

\newcommand{\cufechapter}[2]{%
  \begin{frame}[plain,noframenumbering]
    \begin{tikzpicture}[remember picture,overlay,x=\paperwidth,y=\paperheight]
      \begin{scope}[shift={(current page.south west)}]
        \fill[cufeblue] (0,0) rectangle (1,1);
        \fill[cufedark] (.72,0)--(1,0)--(1,1)--(.90,1)--cycle;
        \node[anchor=west,text=white,font=\fontsize{40}{44}\selectfont\bfseries]
          at (.075,.66) {#1};
        \draw[cufegold,line width=1.2pt] (.08,.52)--(.22,.52);
        \node[anchor=west,text=white,font=\LARGE\bfseries] at (.075,.40) {#2};
        \node[anchor=south east,inner sep=0] at (.93,.08) {\cufeseal{1.1cm}};
      \end{scope}
    \end{tikzpicture}
  \end{frame}
}
\newcommand{\backmatter}{%
  \begin{frame}[plain,noframenumbering]
    \begin{tikzpicture}[remember picture,overlay,x=\paperwidth,y=\paperheight]
      \begin{scope}[shift={(current page.south west)}]
        \fill[white] (0,0) rectangle (1,1);
        \node at (.5,.73) {\cufeseal{1.25cm}};
        \node[text=cufeblue,font=\LARGE\bfseries] at (.5,.48) {感谢聆听，敬请指正};
        \node[text=cufetext,font=\small] at (.5,.32) {\insertauthor\quad |\quad 中央财经大学};
        \node[text=cufegray,font=\small] at (.5,.18) {忠诚\quad 团结\quad 求实\quad 创新};
      \end{scope}
    \end{tikzpicture}
  \end{frame}
}

% ---------- 从这里修改报告信息 ----------
\title[学术报告]{报告标题}
\subtitle{报告副标题}
\author{作者姓名}
\institute{中央财经大学\quad 学院／研究机构}
\date{2026年10月}
% \titlebackground*{assets/campus.jpg} % 如需校园照片，上传图片并取消注释。

\begin{document}
\begin{CJK*}{UTF8}{gbsn}
\maketitle

\begin{frame}{报告目录}
  \tableofcontents
\end{frame}

\section{研究背景}
\cufechapter{01}{研究背景}
\begin{frame}{研究背景与核心问题}
  \framesubtitle{用一页说明研究动机、问题与价值}
  \begin{itemize}
    \item \textbf{研究背景：}简要介绍现实背景和已有研究。
    \item \textbf{核心问题：}明确本报告希望回答的问题。
    \item \textbf{研究价值：}说明理论贡献或实践意义。
  \end{itemize}
  \vspace{3mm}
  \begin{block}{本页要点}
    每页围绕一个核心观点组织内容，标题尽量直接表达结论。
  \end{block}
\end{frame}

\begin{frame}{研究思路与分析框架}
  \begin{columns}[T,onlytextwidth]
    \begin{column}{.46\textwidth}
      \begin{block}{研究设计}
        \begin{itemize}
          \item 识别关键变量
          \item 构建理论框架
          \item 选择合适方法
        \end{itemize}
      \end{block}
    \end{column}
    \begin{column}{.49\textwidth}
      \centering
      \begin{tikzpicture}[
        box/.style={draw=cufeblue,fill=cufelight,text width=3.4cm,
          minimum height=.85cm,align=center,font=\small},
        link/.style={-{Stealth},thick,cufeblue},node distance=5mm]
        \node[box] (a) {研究问题};
        \node[box,below=of a] (b) {理论机制与研究假设};
        \node[box,below=of b] (c) {实证检验与政策启示};
        \draw[link] (a)--(b);
        \draw[link] (b)--(c);
      \end{tikzpicture}
    \end{column}
  \end{columns}
\end{frame}

\section{研究方法与结果}
\cufechapter{02}{研究方法与结果}
\begin{frame}{模型设定与变量说明}
  \begin{block}{示例：面板数据回归模型}
    \[
      y_{it}=\alpha+\beta x_{it}+\boldsymbol{\gamma}^{\mathsf T}
      \mathbf{Z}_{it}+\mu_i+\lambda_t+\varepsilon_{it}
    \]
  \end{block}
  \begin{itemize}
    \item $y_{it}$：被解释变量；$x_{it}$：核心解释变量。
    \item $\mathbf{Z}_{it}$：控制变量；$\mu_i$、$\lambda_t$：个体与时间固定效应。
    \item $\beta$：待估计的核心参数；识别解释需结合研究设计。
  \end{itemize}
\end{frame}

\begin{frame}{主要结果展示}
  \framesubtitle{以下为排版示例，所有数值均为演示数据}
  \centering
  \begin{tabular}{lccc}
    \toprule
    变量 & 模型（1） & 模型（2） & 模型（3） \\
    \midrule
    核心解释变量 & 0.125 & 0.118 & 0.109 \\
    标准误 & (0.042) & (0.039) & (0.037) \\
    控制变量 & 否 & 是 & 是 \\
    个体固定效应 & 否 & 否 & 是 \\
    样本量 & 1,200 & 1,200 & 1,200 \\
    \bottomrule
  \end{tabular}
  \vspace{4mm}
  \begin{block}{结果解读}
    在此解释估计结果、经济含义以及结论成立的条件。
  \end{block}
\end{frame}

\section{结论与讨论}
\cufechapter{03}{结论与讨论}
\begin{frame}{主要结论与后续方向}
  \begin{enumerate}
    \item \textbf{主要发现：}用简明语言回答最初提出的问题。
    \item \textbf{实践启示：}给出与研究证据相匹配的建议。
    \item \textbf{研究局限：}说明数据、方法和适用范围的限制。
  \end{enumerate}
  \vspace{3mm}
  \begin{alertblock}{后续方向}
    在此补充进一步研究的问题或讨论议题。
  \end{alertblock}
\end{frame}

\begin{frame}[fragile]{模板使用说明}
  \begin{itemize}
    \item 默认使用 \textbf{pdfLaTeX}；请一并上传模板包内的 assets 文件夹。
    \item 修改导言区的标题、作者、学院与日期，再替换示例正文。
    \item 使用 \verb|\cufechapter{01}{章节标题}| 插入章节页。
    \item 使用 \verb|\footlinecolor{cufeblue}| 调整页脚配色。
    \item 本模板为个人制作的中财风格版本，非学校官方发布。
  \end{itemize}
  \vspace{2mm}
  {\scriptsize 版式参考：\href{https://www.overleaf.com/latex/templates/bei-da-zhong-wen-mo-ban-pku-beamer-template/kfxpbtzrqhrn}{PKU-Beamer-Template}，Zhuming Shi，CC BY 4.0。\\
  原始主题来源：Federico Zenith（SINTEF Presentation）、Liu Qilong。\\
  校徽使用中央财经大学官网公开标识；本版配色为屏幕设计用色。}
\end{frame}

\backmatter
\end{CJK*}
\end{document}
